\documentclass[10pt,a4paper]{amsart}
\usepackage[margin=19mm]{geometry}
\usepackage{amsmath,amssymb,mathtools}
\usepackage[scr=rsfs,cal=euler]{mathalfa}
\usepackage{xfrac}
\usepackage[unicode,hidelinks,bookmarks=false]{hyperref}
\newcommand{\ie}{i.e.\ }
\newcommand{\cf}{cf.\ }
\newcommand{\resp}{respectively\ }
\newcommand{\Subsection}{\subsection}
\providecommand{\nobreakdash}{\nobreak\hbox{-}}
\setlength{\emergencystretch}{2em}
\multlinegap0pt
\usepackage{bm}
\usepackage{bbm}
\usepackage{dsfont}
\usepackage{xspace}
\usepackage{cancel}
\usepackage{mathtools}
\usepackage{amsthm}
\makeatletter
\@ifundefined{theorem}{\newtheorem{theorem}{Theorem}}{}
\@ifundefined{lemma}{\newtheorem{lemma}[theorem]{Lemma}}{}
\makeatother
\makeatletter
\newcommand{\floor}[1]{\left\lfloor#1\right \rfloor}
\expandafter\ifx\csname pfofthm\endcsname\relax
\newenvironment{pfofthm}[1]
{\par\vskip2\parsep\noindent{\sc Proof of\ #1. }}{{\hfill
$\Box$}
\par\vskip2\parsep}
\fi
\makeatother
\title[Two-stage Bootstrap Percolation]{Two-stage Bootstrap Percolation}
\author{Zihao Fang, Janko Gravner, David Sivakoff}
\date{Source: arXiv:2509.16541v1 (CC BY 4.0)}
\begin{document}
\maketitle

\noindent\emph{Source and licence.} Two-stage Bootstrap Percolation. Version arXiv:2509.16541v1 (CC BY 4.0), distributed under the Creative Commons Attribution 4.0 International licence, as recorded in the arXiv licence metadata for this exact version and shown on the abstract page https://arxiv.org/abs/2509.16541v1. Full rights notice: licenses/arxiv-2509.16541-CC-BY-4.0.txt.\par\smallskip

\noindent\textit{Editorial scope.} Counted unit: the Lemma whose source label is \texttt{lem:F3} in the article (source0.tex lines 1172--1174, source label \texttt{lem:F3}), delivered together with the article's own complete original proof of it (source0.tex lines 1175--1249, one \texttt{proof} environment of 75 lines). No mathematics is written, repaired or re-derived: statement and proof are the article's own text line for line. Required context retained uncounted: none. Every definition, display and label of those lines is retained unchanged. Omitted from this package and not redistributed: 1--1171, 1250--2990. Nothing inside a retained excerpt is omitted or altered: every retained body is delivered exactly as the article's lines are written, so no omission receipt of any kind accompanies this package. Citations: the retained excerpts carry no citation command at all, so no bibliography entry of the article is transcribed and no reference is redistributed. Reference adaptation: none was needed and none was made; every reference inside the delivered unit resolves inside it, so no unresolved reference or citation is produced. Printed slips kept literal: no printed word, symbol, hypothesis or number of the counted statement or of its proof was altered. Layout and class: the article's own class is \texttt{article} with packages including babel, graphicx, amssymb, amsmath, amsthm, fullpage, mathrsfs, subfigure, verbatim, enumitem, bbm, tikz, authblk, indentfirst; the wrapper is the lane's pinned \texttt{amsart} editorial wrapper at 10pt on A4, which restates the theorem-like environments the retained text needs and adds the article's own macro definitions that the retained text needs (\texttt{\textbackslash floor}), with extra wrapper packages (bm, bbm, dsfont, xspace, cancel, mathtools). The delivered numbering is the unit's own. Assets: no retained span contains a figure environment, a graphics-inclusion command, a PDF-page inclusion, a table or a TikZ picture; no image file is copied into this package, the package declares no asset dependency and no image file is redistributed with it. Licence: version arXiv:2509.16541v1 of this article is distributed under the Creative Commons Attribution 4.0 International licence, as recorded in the arXiv licence metadata for this exact version and shown on the abstract page https://arxiv.org/abs/2509.16541v1. A full rights notice is delivered at licenses/arxiv-2509.16541-CC-BY-4.0.txt. Characters: every retained excerpt is pure ASCII, so the delivered engine, XeLaTeX with its default Latin Modern text font, reports no missing character.
\begin{lemma}\label{lem:F3}
    A $q$-box satisfies condition (F3) with high probability provided that $q \le p^2$.
\end{lemma}

\begin{proof}
    First, the expected number of $5\times 5$ squares within the $q$-box that contain at least three initial $2$s is at most a constant times 
    \[
    \left((1/q)\log(1/q)\right)^2 \cdot q^3
    = q\cdot (\log(1/q))^2\to 0.
    \]
    

    Now cover the $q$-box by squares of side length $6 \floor{\frac{1}{p}\log(1/p)\log(1/q)}$, which half overlap with neighboring squares. 
    The number of such squares required to cover the $q$-box is at most
    a constant times 
    \begin{equation}\label{eqn:F3-squares}
    \left(\frac{(1/q)\log(1/q)}{(1/p)\log(1/p)\log(1/q)}\right)^2 = \frac{1}{(\log(1/p))^2} \cdot \frac{p^2}{q^2}.
    \end{equation}
    
    % If a square with side length $3 \floor{\frac{1}{p}\log(1/p)\log(1/q)}$ in the $q$-box initially contains a frame, then one of these larger squares initially contains a frame. 
    % By definition, a frame requires at least $3$ and at most $8$ initial state-$2$ vertices.
    % Notice that a frame defined by only three initial state-$2$ vertices must be contained in a $5 \times 5$ square. \note{JG: as two sides at distance at least 5 cannot share the requisite state-$2$ vertices.}
    % So by the first part of the proof, such a frame defined by exactly three $2$s do not exist in the $q$-box with high probability.
    % For a frame defined by exactly four $2$s, all these four state-$2$ vertices must be within $\ell^\infty$-distance $2$ to the corners of the frame.
    % Then with $q \leq p^2$, the expected number of the larger squares that contain a frame defined by exactly four $2$s is at most
    % \[
    % \cO\left(\frac{1}{(\log(1/p))^2} \cdot \frac{p^2}{q^2} \cdot \left(\frac{6\log(1/p)\log(1/q)}{p} \right)^4 \cdot q^4 \right)
    % = \cO\left(q (\log(1/q))^6\right) \to 0
    % \]
    % as $q \to 0$. 

    % When a frame is defined by exactly five $2$s, three of them must be placed around a pair of opposite corners (two $2$s around one corner and the third around the other corner), and the remaining two $2$s can be freely placed along the two sides connected to the corner with only one $2$.
    % The expected number of larger squares that contain such a frame is at most
    % \[
    % \cO\left(\frac{1}{(\log(1/p))^2} \cdot \frac{p^2}{q^2} \cdot \left(\frac{6\log(1/p)\log(1/q)}{p} \right)^4 \cdot \left(\frac{6\log(1/p)\log(1/q)}{p} \right)^2 \cdot q^5 \right)
    % = \cO\left(q(\log(1/q))^{10}\right) \to 0
    % \]
    % as $q \to 0$. 

    % For a frame defined by exactly six $2$s, only two $2$s need to be around one corner, and the remaining four $2$s have freedom along the two sides connected to the opposite corner.
    % The expected number is then at most
    % \[
    % \cO\left(\frac{1}{(\log(1/p))^2} \cdot \frac{p^2}{q^2} \cdot \left(\frac{6\log(1/p)\log(1/q)}{p} \right)^4 \cdot \left(\frac{6\log(1/p)\log(1/q)}{p} \right)^4 \cdot q^6 \right) \to 0
    % \]
    % as $q \to 0$.

    % With exactly seven $2$s, only one $2$ needs to be around one corner to be shared, and all the other six $2$s have freedom.
    % The expected number is at most
    % \[
    % \cO\left(\frac{1}{(\log(1/p))^2} \cdot \frac{p^2}{q^2} \cdot \left(\frac{6\log(1/p)\log(1/q)}{p} \right)^4 \cdot \left(\frac{6\log(1/p)\log(1/q)}{p} \right)^6 \cdot q^7 \right) \to 0
    % \]
    % as $q \to 0$.

    % Finally, with eight $2$s, every $2$ has freedom along the sides and no $2$ needs to be shared.
    % The expected number is at most
    % \[
    % \cO\left(\frac{1}{(\log(1/p))^2} \cdot \frac{p^2}{q^2} \cdot \left(\frac{6\log(1/p)\log(1/q)}{p} \right)^4 \cdot \left(\frac{6\log(1/p)\log(1/q)}{p} \right)^8 \cdot q^8 \right) \to 0
    % \]
    % as $q \to 0$.

    If the longest side of a frame has at least 6 sites, then there are two parallel lines at distance at least 5, each with two state-$2$ sites within distance 2 of it. In addition, there is a perpendicular line that has within distance 2 either: 
    two of these four $2$s;
    or one of these four 2s and an additional $2$; or two additional $2$s. (If the first two lines are vertical, then the third line may be, say, the top line of the frame, which must have two $2$s within distance $2$.) In order, the three scenarios have their probabilities bounded by $(\log(1/p)\log(1/q))^{10}$ times: 
    \begin{equation*}
        \begin{aligned}
         &\frac{p^2}{q^2} \cdot \frac{1}{p^2} \cdot  \frac{q^4}{p^3}=
    \frac{q^2}{p^3};\\
    &
    \frac{p^2}{q^2} \cdot \frac{1}{p^2} \cdot  \frac{q^5}{p^5}=
    \frac{q^3}{p^5}; \text{ and}
    \\&
    \frac{p^2}{q^2} \cdot \frac{1}{p^3} \cdot  \frac{q^6}{p^6}=
    \frac{q^4}{p^7}.   
        \end{aligned}
    \end{equation*}
   For example, the first factor in the second probability bound 
   represents the number of squares needed to cover the $q$-box bounded by (\ref{eqn:F3-squares}) and the second factor bounds the number of lines within a fixed square. Finally, for the third factor, we first choose the four $2$s on the two original lines, which gives the factor $(q/p)^4$, and then an additional $2$ on a perpendicular line through one of the chosen four locations, which gives the additional factor $q/p$. 
    As $q\le p^2$, the three probabilities approach $0$ as $p\to 0$.   
\end{proof}

\end{document}
